\documentclass[11pt]{article}
\usepackage{graphicx}
\usepackage[margin=1in]{geometry}
\usepackage{amsmath,amssymb,amsthm}
\usepackage{booktabs}
\usepackage{hyperref}
\usepackage[numbers]{natbib}
\newtheorem{theorem}{Theorem}
\newtheorem{lemma}[theorem]{Lemma}
\newtheorem{definition}[theorem]{Definition}
\newtheorem{proposition}[theorem]{Proposition}
\newtheorem{corollary}[theorem]{Corollary}
\theoremstyle{remark}
\newtheorem{remark}[theorem]{Remark}

\title{One bit stops the walk: cycle detection in a two-header tape machine}
\author{HaoWei Luo\thanks{ORCID: 0009-0007-7206-7095; \texttt{luohaowei@bjtu.edu.cn}.}}
\date{November 2026}

\begin{document}
\maketitle

\begin{abstract}
A tape machine with two headers, a window and three trims answers a query at each step; in the companion
paper its cost is measured on the tapes whose walk \emph{meets}. This paper asks how many of them do, and
what the least possible memory is that turns the rest into runs with an answer. Over $18{,}000$ measured
(rule, $n$, tape) cells, $14.9\%$ of the runs meet and $61.9\%$ close a period and then repeat it for ever,
dying at the step cap; the cycle is short, a median of two steps. One bit --- the answer to ``has this
window occurred before?'' --- takes the meeting rate to $37.8\%$ and rescues $36.9\%$ of the runs that
otherwise die. The bit is not state: the visited set is the oracle's, the machine keeps $O(\log n)$ bits.
It buys termination and no cheapness: a run that used to die now continues to an answer, trims cascade (a
median rescued tape needs a few hundred, the worst seen $1{,}956$), and the read bound is
$O(n^{3})$, no better than the most expensive rule in the catalogue. We also record why one bit of \emph{ink} --- a writable
mark per cell --- buys nothing at all: a mark read costs what a cell read costs, so a mark can only pay by
answering a question about several cells at once, and one bit cannot.
\end{abstract}

\section{The problem}

The companion paper \citep{luo-bthm-v2} studies a tape machine whose walk is a pair of headers enclosing a
window on a tape of $n$ three-valued cells. At each step a fixed rule, given the window's content and the
geometry, names one of eight actions: move a header inward or outward, trim a cell from the left, the right
or both, or halt. A tape \emph{meets} when the two headers coincide; only a meeting tape carries an answer
$\delta$, and every cost in that paper is conditional on meeting.

The condition is not a formality. Measuring the walk of twenty rules on the standard ensemble
($300$ tapes per cell, $n=512,1024,2048$; Section~\ref{sec:measure}) gives $18{,}000$ tape-runs in sixty $(rule,n)$ cells, of which
$2{,}681$ ($14.9\%$) meet and $11{,}146$ ($61.9\%$) close a period and then repeat it for ever, hitting the
step cap without an answer. For most rules the meeting rate at $n=512$ lies between $2.3\%$ and $31.7\%$
(Table~\ref{tab:cyc}); only \texttt{monotone} and \texttt{hashmod} meet on every tape. The companion paper's means
are therefore taken over a small and rule-dependent subpopulation, and its tables do not report the share
(its appendix gives pal's meeting rate, 43 of 300 draws at $n=512$, but no other rule's).

This paper asks the corresponding question of memory. A walk that repeats a state repeats it for ever,
because the next state is a function of the state alone; the walk cannot tell that it has been here before,
and no amount of time helps. What is the least memory that lets it tell? We give the machine exactly one
bit per state --- an oracle answering \emph{has this window occurred before in this run?} --- and measure
what that buys. The answer has two halves, and the second is negative: it buys termination, and it buys no
cheapness.

\section{The model}

\paragraph{The walk.} The companion paper's model, unchanged. A region $[L,R]$, a window
$[\ell,r]\subseteq[L,R]$, the state $\sigma=(\ell,r,L,R)$; eight actions, three of which trim a cell and
reset both headers; a tape meets when $\ell=r$. The rule is a function $f(\varphi(W),n,\ell,r,L,R)$ with
$\varphi$ a fixed function of the window's content.

\paragraph{The bit.} Add one query, available at every state:
\[
\operatorname{seen}(\ell,r,L,R)=1 \iff \text{this quadruple has occurred earlier in the same run.}
\]
A rule may branch on it; on a repeat it may trim and continue. Nothing else is added. The visited set is the
oracle's, not the machine's: the machine's own state remains $O(\log n)$ bits, and it cannot store, count or
name the states it has seen. In particular the oracle does not answer \emph{how many times} a window has
occurred, and the rule cannot recover a count from it (Section~\ref{sec:not}).

\paragraph{The trim side.} The bit says that the state has occurred; it does not say which end to trim. This
is a real degree of freedom: trimming the left end and trimming the right differ by a median of $0.010$ in
meeting fraction and by at most $0.107$ (\texttt{AMAX}, \texttt{AMIN}). The model therefore declares the side rather
than leaving it to the implementation: on a repeat the rule trims the side its own table selects for that
state --- the same parity rule it uses when it cuts --- one cell per repeat. The repair is a trim in the
companion paper's sense: it removes the boundary cell of the \emph{tape}, so the region shrinks. It is not a
withdrawal of the window, which the walk could undo by widening it again; the difference matters in
Section~\ref{sec:term}.

\begin{lemma}[Containment]\label{lem:contain}
Every rule of the companion paper is a rule of this model with the same cost on every tape. A rule that
halts without an answer when the bit is $1$ behaves exactly as in the companion paper: the tapes on which it
halts without an answer are the tapes the companion paper does not count either.
\end{lemma}
\begin{proof}
A rule that never reads the bit is a companion-paper rule. For the second sentence, the run of such a rule
is the companion-paper run up to the first repeat; at the repeat it halts, and the tape is one on which no
answer is produced, which is precisely a tape the companion paper excludes. The same holds for a rule that
trims on the bit rather than halting: a repeated quadruple means the walk is periodic from there on, so such
a tape never meets in the companion paper either, and the bit is silent on every tape both models count.
\end{proof}

So the bit cannot change the cost of a tape both models count (Lemma~\ref{lem:contain}). It can only change the \emph{domain}: the set
of tapes on which an answer is produced. That is the whole of its effect, and Section~\ref{sec:measure}
measures it.

\section{Termination, and what it costs}\label{sec:term}

\begin{theorem}\label{thm:term}
If a rule trims a tape cell whenever the bit is $1$, every run halts in at most $n+n(n^{2}+1)$ steps.
\end{theorem}
\begin{proof}
Every trim removes a cell of the tape --- the rule's own cuts and the repair alike --- so there are at
most $n$ trims in the whole run. Between two of them the region $[L,R]$ is fixed, so the state is the pair
$(\ell,r)$ with $L\le\ell\le r\le R$: at most $n^{2}$ states. If the walk returned to a state it had already visited, the
bit would be $1$ there and the rule would trim, so no state repeats between trims; hence at most $n^{2}$
steps between trims, and with at most $n$ trims the run has at most $n+n(n^{2}+1)$ steps.
\end{proof}

\begin{corollary}\label{cor:cost}
For the collapse family, where a step costs one read and a trim one more, the read cost of a terminating run
is $O(n^{3})$.
\end{corollary}

The bound of Theorem~\ref{thm:term}, and of Corollary~\ref{cor:cost} for the reads, is not an
improvement: it is no better than the most expensive rule in the published catalogue, whose cheapest
family costs $n+S$ and whose most expensive costs $n^{3}$. The bit guarantees that a run ends, not that it
ends cheaply, and two features make that concrete. First, the repair itself costs: a median rescued tape
needs a few hundred trims and the worst needs $1{,}956$ (\texttt{NBLK=1}, $n=2048$), each paid at one read
plus the head re-read. Second, the rescued runs are the expensive ones: a tape that used to die at the cap
now continues to an answer and pays for the continuation, so the median reads over all tapes rise even as
the meeting rate does.

\section{The measurement}\label{sec:measure}

\paragraph{Harness.} The ensemble of the companion paper: $300$ random tapes per cell, $n=512,1024,2048$,
one charge for every rule so that the two models are comparable (one read per cell read, counted with
multiplicity; the oracle bit is free; a trim pays the cell it gives up plus the head re-read). A cycle is a
repeated quadruple; since the tape is fixed, the quadruple determines the next state, so a repeat is a closed
period and the run then dies at the cap ($40n^{2}+100$ steps). The repair is the trim of
Section~\ref{sec:term}: on a repeat the rule removes the boundary cell of the tape, on the side its own table
selects, and continues.

\paragraph{The gates.} Before any number was recorded: $99/99$ cells exact against the witness panel's cells
at $n=512$ (reads, $\delta$, end); $114/114$ $(\delta,\text{end})$ against the reference engines on six
random tapes and twenty rules; $800/800$ detected periods re-run for a full period; $30/30$ tape-runs
matching a second, independently written engine on all six fields of the repair; and, on every one of the
$18{,}000$ tape-runs, the no-memory walk reproduced field by field against the earlier run ($0$ mismatches on
seven fields), which is what makes the two runs comparable. The repair never fires more than $n$ times
(Table~\ref{tab:agg}), no run hit the step cap, and no cell ran out of oracle table.

\begin{table}[t]
\centering
\small
\input{probe_cycle_cut/_cut_table.tex}
\caption{One row per rule at $n=512$, $300$ tapes, with the cut repair. \emph{BTHM-0}: the share of tapes
that meet, and the share that close a period and die at the cap. \emph{cycle}: the median period length in
steps and the median window width at the repeat. \emph{BTHM-C}: the meeting share, the rescued share of the
cap-diers, and the median number of repair trims a rescued tape needs. \emph{reads}: the median over the
tapes the companion paper counts, and over all tapes under the bit.}
\label{tab:cyc}
\end{table}

\paragraph{What the bit buys.} Over the $18{,}000$ tape-runs the meeting rate goes from $14.9\%$ to $37.8\%$,
and $4{,}115$ of the $11{,}146$ cap-diers are rescued ($36.9\%$). No run under the bit hit the cap. The
rescue is uneven: it is $1.000$ for \texttt{hashmod} and \texttt{monotone}, which never cycle, and between
$0.28$ and $0.50$ for the rules that do, the largest for \texttt{dense2} ($0.503$ at $n=512$) and the
smallest for \texttt{runs} ($0.276$).

\paragraph{The shape of a cycle.} The walk does not wander before it repeats: the median period is
\textbf{two steps}, a move out and the same move back, at every size and for most rules. The width at which
the repeat happens is rule-specific (Figure~\ref{fig:width}) --- five cells for \texttt{pal}, eleven for
\texttt{longrun}, about a third of the tape for \texttt{NBLK>2} and \texttt{BAL>m/2} ($749$ and $720$ cells
at $n=2048$). A long period is the exception: \texttt{dense2} cycles in $13$--$14$ steps, \texttt{pattern}
in six.

\begin{figure}[t]
\centering
\includegraphics[width=0.80\textwidth]{probe_cycle/figs/cyc_width.pdf}
\caption{Where the walk repeats itself, at $n=512$: the median window width at the repeat as a fraction of
the tape, with the median period written beside each bar. Four rules repeat only after the window has
reached a third of the tape or more, and those are the rules whose repairs are most expensive.}
\label{fig:width}
\end{figure}

\paragraph{The two rules that never need it.} \texttt{monotone} has no trim and increases $\ell$ every step,
so no quadruple can repeat; it meets all $300$ tapes in $n-1$ steps at $2n-1$ reads, and the bit is worth
exactly nothing. \texttt{hashmod} meets all $300$ tapes in both models, so again nothing.

\subsection{The repair, and what it decides}\label{sec:repair}

The repair is where the model is tested, in three ways.

\emph{It decides whether there is a theorem.} The termination proof of Section~\ref{sec:term} needs every
repair to consume a cell of the tape. A cheaper repair suggests itself: retract the window one cell and leave
the tape untouched. Measured, that repair is far better at rescuing runs --- $74.0\%$ of them meet against
$37.8\%$ --- and it hits the cap on no tape either; but the walk can undo it by re-widening the window, so its
repair count has no bound ($3{,}627{,}840$ trims at $n=2048$, against $1{,}956$ for the cut) and no
termination theorem applies to it. This paper reports the cut, because a model whose termination is observed
but not proved is a weaker claim than one whose termination is proved.

\emph{It decides the rescue rate, by a factor of two.} Table~\ref{tab:repair} gives the two repairs side by
side at $n=512$: the retraction rescues up to three times as many runs on the rules that cycle most
(\texttt{peak} $0.813$ against $0.340$, \texttt{BAL>m/2} $0.887$ against $0.393$), and the gap is a property
of the repair, not of the bit.

\emph{Not every repair can work.} A repair that only reverses the last move does not terminate: it returns
to a state that has also been visited, so the bit fires again and the rule reverses again, and the walk
alternates between two seen states. Escaping needs a change in the geometry, and a reversal changes none.

\begin{table}[t]
\centering
\small
\input{probe_cycle_cut/_cutagg_table.tex}
\caption{The aggregate at the three sizes, over the twenty rules that were run (\texttt{monotone} is
analytic and enters no cell). The meeting rate falls slowly with $n$ while the cap share rises: divergence is
not a small-size artefact. The last column is the largest number of repair trims a single run needed; it is
always at most $n$, as the theorem requires.}
\label{tab:agg}
\end{table}

\begin{table}[t]
\centering
\small
\input{probe_cycle_cut/_cutrep_table.tex}
\caption{The repair decides the rescue rate: the window retraction against the cut, at $n=512$, as the share
of tapes that meet. The retraction is better everywhere and far better on the rules that cycle most, and it
is the repair for which no termination bound holds.}
\label{tab:repair}
\end{table}

\subsection{What the bit costs, and why it is not state}\label{sec:account}

Three accountings could hide in the bit, and the paper fixes all three.

\emph{It is not stored state.} The set of visited quadruples has size $\Theta(n^{4})$, so no machine with
$O(\log n)$ bits can hold it: the oracle answers about a set the machine could never carry. The model is
therefore not the companion paper with a larger state budget in disguise.

\emph{It costs no reads.} The query is free. What the bit changes is the number and the length of runs, and
the price is paid in trims: each costs one read for the cell it gives up plus the head re-read, and the trim
counts are in Table~\ref{tab:cyc}.

\emph{It does not count.} The oracle answers a Boolean, so a rule cannot recover a visit count, its parity,
or a colour from it. The counting variant is a different model, and Section~\ref{sec:not} says why.

\section{Why an oracle, and not ink}\label{sec:not}

The obvious alternative is to give the machine a writable mark per cell instead of an oracle: one bit beside
each cell's data, readable with it, writable at the cost of one write. It buys nothing, and the reason is
structural rather than a matter of the rules tried.

A read returns the cell's data and its mark together for one read, so consulting a mark costs exactly what
consulting the cell costs. A mark can therefore pay only by answering a question about \emph{several} cells
at once. One bit can do that in two ways. It can encode a fixed-offset relation, which pays only where the
charge is already per pair --- and \texttt{pal} is already that case, one comparison being one read for two
cells. Or it can carry a monotone Boolean range summary, which pays only where the predicate is not already
free --- and \texttt{pattern}'s predicate is free in its engine. Measured on the two rules for which the record
has real machines: both exact variants of \texttt{pal} read exactly the baseline at every size (ratio $1.0000$)
and only spend ink; its one worthwhile certificate, the diagonal, needs about $27$ bits at $n=512$ and recurs
a median of $129$ steps later, so an untagged bit is stale exactly when it would be useful; a position-only
version of the same idea reads $12.6$--$21.6$ times more and returns the wrong answer on $91/85/98$ of $300$
tapes. \texttt{pattern} has one working mechanism, an adjacency bit on the paired opening pass, worth
$190/383/768$ reads --- $0.25\%$ of the charge falling to $0.07\%$ --- fitted exponent $1.009$ against a
charge that fits $1.963$: it does not scale. Even an unbounded, perfectly tagged certificate store would save
$17.1/18.0/18.6\%$ of the charge, and the plane holds about $18$ of the roughly $5{,}740$ certificates one
tape needs.

A larger plane is not a larger version of this model; it is a different one, and its question is different:
what does $O(n)$ bits of state, on the tape, buy over $O(\log n)$?
That question is not answered here.

\section{What this says about the companion tables}

Two consequences, and the second is a correction of interpretation rather than of number.

First, the companion paper's costs are not the costs of the model on all tapes; they are the costs on the
$14.9\%$ that meet, and the rate varies by a factor of forty across rules (from $2.3\%$ for the level-count
rules to $100\%$ for \texttt{hashmod}). Any comparison between two rules' means is therefore a comparison of two
different subpopulations, and a revision should print the rate beside the mean.

Second, divergence is not a defect of particular rules: it is what the model does. The walk is memoryless,
so a rule that can return to a state will return to it, and the state space is $O(n^{4})$. Of the twenty
rules measured, eighteen cycle on some tapes, and for fifteen of them the cycle is the majority outcome at some
size. What saves the meeting tapes is geometry, not design: the walks that shrink monotonically cannot
repeat, and the walks that end quickly escape the cap.

\section{Open}

\paragraph{Can a machine hold the bit?} The oracle answers about a set of size $n^{4}$; no machine of the
class can store it. For a monotone walk the window is determined by the step count, so the bit is free and
the walk never needs it. For a general walk the question is whether some rules can recognise their own
repeats with $O(\log n)$ bits --- a question about the rule, not about the model, and one this paper does not
answer.

\paragraph{Counting instead of remembering.} An oracle that answers \emph{how many times} (or the count
modulo $k$) is a different model: its memory is $O(n^{2}\log k)$ bits, it is Langton's-ant-shaped
\citep{langton1986}, and the interesting question there is not cost but computational power. It should not
share a paper with this one.

\paragraph{The domain, as a quantity.} The companion paper reports a cost; this paper reports a domain. A
model that reports both --- the cost on the tapes it answers, and the share of tapes it answers --- would be
strictly more informative than either, and the tables here are a first instance of the second column.

\appendix
\section{Engines and gates}

The walk is implemented once, as a table-driven machine (the action table and the two predicates
$b_1,b_2$ of each rule, taken from the reference engines), and charged with one accounting for every rule so
that the two models are comparable; the reference engines charge differently (band, \texttt{palmach2},
\texttt{hk\_engine}), and mixing them would confound the model with the engine. \texttt{hashmod} is the one rule
whose reference machine is not the common engine; it is nevertheless the same walk, and its $(\delta,\text
{end})$ agrees with the table-driven engine on all six check tapes and its meeting count on all $300$ tapes
at every size. \texttt{runs}, \texttt{pal}, \texttt{longrun} and the band rows were checked the same way.

\begin{thebibliography}{9}
\bibitem[Luo(2026)]{luo-bthm-v2} H.~Luo. A memoryless two-header tape machine: what a window costs, and why
two lower-bound techniques fail. Preprint, version 2.0, \texttt{doi:10.5281/zenodo.24199236}, 2026.
\bibitem[Langton(1986)]{langton1986} C.~G. Langton. Studying artificial life with cellular automata.
\emph{Physica D} 22 (1986), 120--149.
\end{thebibliography}

\end{document}
